\chapter*{Jak używać tej książki}\addcontentsline{toc}{chapter}{Jak używać tej książki}\markboth{Jak używać tej książki}{Jak używać tej książki}

Książka jest zaplanowana na dziesięć tygodni, po jednym rozdziale na tydzień. Ten rytm nie jest obietnicą. Rozdział o pracy może zająć dwa tygodnie, a rozdział o chlebie jeden wieczór. Ważniejsza od tempa jest kolejność. Najpierw liczysz dobę (Część I), potem zmieniasz rytm dnia (Część II), a dopiero potem wchodzisz w pieczenie (Część III). Jeśli przeskoczysz od razu do chleba, dostaniesz receptury bez planu, w którym mogą się zmieścić.

\section*{Co robić w każdym rozdziale}

Zacznij od celów. Są krótkie i mówią, co powinieneś umieć po rozdziale. Przy każdym przykładzie policz sam, zanim spojrzysz na wynik. Ramki mają konkretne role. \emph{Pułapka} opisuje błąd, w który wpada wiele osób, i jest tam po to, żebyś najpierw sam na niego wpadł, a dopiero potem przeczytał korektę. \emph{Zasada} to reguła wartą zapamiętania. \emph{Uwaga} to zastrzeżenie, którego nie wolno pominąć. \emph{Przykład} to rozpisane liczby, na których widać mechanizm. Ramka \emph{Z piekarnika} zbiera uwagi dotyczące pieca i zawiera tylko te, których nie zastąpi zdrowy rozsądek.

\section*{Ćwiczenia i zadanie na tydzień}

Ćwiczenia są na końcu rozdziału. Odpowiedzi zebrano w Dodatku A, ale zanim do nich zajrzysz, spróbuj każdego ćwiczenia na kartce. Odpowiedź, którą przeczytasz przed próbą, nie uczy niczego. Liczba ćwiczeń w książce jest liczona przez narzędzia, a nie wpisywana ręcznie, i zmienia się wraz z rozdziałami.

Zadanie na ten tydzień jest jedyną rzeczą w rozdziale, która ma zostać wykonana, a nie tylko przeczytana. Ma miarę sukcesu. Wynik zapisz w swoim szablonie z Dodatku C, bo tylko wtedy da się go porównać z następnym tygodniem.

\section*{Oznaczenia}

Pomarańczowy znaczek \textsuperscript{\textcolor{accentorange}{[?]}} (w indeksie górnym, w nawiasach kwadratowych) oznacza twierdzenie, którego autor nie sprawdził w źródle pierwotnym. To nie znaczy, że twierdzenie jest błędne. Znaczy, że nie powinieneś na nim budować decyzji, zanim go nie sprawdzisz. Jak to zrobić, opisuje Dodatek D, a liczbę takich miejsc podaje rejestr w Dodatku Z. Liczby bez znaczka są albo policzone z danych w książce, albo są arytmetyką pokazaną w tekście, którą możesz powtórzyć.

\section*{Tryb minimum}

Będą tygodnie, w których nie uda się nic z planu. Wtedy obowiązuje tryb minimum, a nie zawieszenie lektury. Składa się z trzech rzeczy. Sen trzymasz w zaplanowanych godzinach. Z szablonu dnia zostawiasz tylko blok snu i blok wolnego. Chleb pieczesz raz w tygodniu, według receptury z najkrótszą listą kroków w Dodatku B. Tryb minimum nie ma wyniku do osiągnięcia. Ma sprawić, że tydzień się nie rozpadnie i że w kolejnym wrócisz do planu bez poczucia porażki.

Jeśli trzy tygodnie z rzędu korzystasz wyłącznie z trybu minimum, to jest informacja o planie, a nie o Tobie. Zmień szablon, zanim zmienisz siebie.
